% ============================================================================
% Appendix: Phase 2 N2 — Bulk Gravitational Backreaction Test
% ============================================================================
% Status: [Dc|model] derivation/bounding attempt within declared 5D model
% Scope: Test whether the bulk gravitational response to junction displacement
%        generates a non-geometric contribution V_bulk(q) to V(q).
% Phase 2 / N2 WP2 deliverable
% Governing documents: PHASE2_NEXTSTEP_PLAN_V1.md,
%                      audit/N2_WP1_DONOR_NORMALIZATION.md
% ============================================================================

\chapter{Bulk Gravitational Backreaction Test for the Non-Geometric Sector}
\label{app:P2_N2_bulk}

% ============================================================================
\section{Scope and Framing}
\label{sec:P2N2:scope}
% ============================================================================

This appendix tests the \textbf{N2 bulk gravitational backreaction lane}
as a source of non-geometric structure in the effective potential $V(q)$.
It builds on the accepted geometric single-well baseline
$V_{\mathrm{geom}}(q) = \tau L_{\mathrm{tot}}(q)$ \tagDc\
(Appendix~\ref{app:vq_chosen_path}) and follows the N2~WP1 donor
normalization and model boundary
(\texttt{audit/N2\_WP1\_DONOR\_NORMALIZATION.md}).

\textbf{What this appendix tests:}
Does the bulk gravitational field response to a displaced Y-junction
produce a $q$-dependent energy contribution $V_{\mathrm{bulk}}(q)$
that can oppose the geometric restoring force and generate new
structure in the combined potential?

\textbf{What this appendix does not claim:}
\begin{itemize}
    \item It does not claim full closure of $V(q)$.
    \item It does not resolve $M_5$, $\Lambda_5$, or the background
          metric choice.
    \item It does not derive $V_B$ or claim double-well existence.
    \item A no-go or bounded insufficiency is an allowed outcome.
\end{itemize}

% ============================================================================
\section{Donor Discipline}
\label{sec:P2N2:donors}
% ============================================================================

\textbf{Allowed donor bundle} (from N2~WP1 normalization):
\begin{itemize}
    \item Put~C corridor structure (D-N2-1) \tagDef\ structural
    \item $V_{\mathrm{geom}}(q) = \tau L_{\mathrm{tot}}(q)$ (D-N2-2) \tagDc
    \item Variant~1--2 negative results (D-N2-3) \tagDc
    \item $\sigma = 8.82\MeV/\fm^2$ \tagDc
    \item $\kappa_5^2 = 8\pi G_5 = 8\pi/\Mfive^3$ \tagDef
    \item $\Mfive$ \tagP\ (not derived; treated parametrically)
    \item WP2 Israel no-go (F-N2-1) \tagDc\ constraint
    \item N7 bounded insufficiency (F-N2-2) \tagDc\ constraint
\end{itemize}

\textbf{Forbidden imports:}
\begin{itemize}
    \item Phenomenological node-well shapes (Variant~3 Gaussian)
    \item $V_B = 2\Delta m_{np}$ as derivation constraint
    \item $\tau_n \approx 878\second$ as input
    \item Background metric chosen post-hoc to produce double-well
    \item $F_{\mathrm{bulk}}$ numerological formula for $G$ (different quantity)
    \item Variant~2 warped scan as positive evidence (it is a negative result)
    \item $C = 100$ or $E_0 = \sigma L_0^2$ (N7 lane quantities)
\end{itemize}

\textbf{Anti-smuggling test} (from N2~WP1 \S7.3):
\emph{Delete all knowledge of $\tau_n$, $V_B$, the Variant~3 form,
and the desired double-well shape. Does the linearized Einstein
equation still yield the same $V_{\mathrm{bulk}}$?}

% ============================================================================
\section{Model-Class Setup}
\label{sec:P2N2:model}
% ============================================================================

\subsection{Background Metric}

\begin{assumptionbox}[title=N2 Model Assumptions (Declared Before Computation)]
\begin{enumerate}
    \item \textbf{Background metric} \tagP:
    \begin{equation}
        ds^2 = e^{2A(\xi)}\,\eta_{\mu\nu}\,dx^\mu dx^\nu + d\xi^2
        \label{eq:N2_background}
    \end{equation}
    Two cases tested: (a)~flat ($A = 0$, $\Lambda_5 = 0$),
    (b)~RS-like ($A = -k|\xi|$, $\Lambda_5 = -6k^2/\kappa_5^2$).

    \item \textbf{5D Planck mass $\Mfive$} \tagP: not derived.
    All results reported as functions of $\Mfive$.

    \item \textbf{Compact direction}: $\xi \in [-\pi R_5, \pi R_5]$
    with $\mathbb{Z}_2$ orbifold \tagP.

    \item \textbf{Junction source}: thin-brane (distributional)
    stress-energy on each arm worldsheet \tagDef.

    \item \textbf{Collective coordinate $q$}: in-brane node
    displacement (inherited from R3) \tagDc.

    \item \textbf{Linearization}: metric perturbation
    $h_{AB}(q) \ll \bar{g}_{AB}$ \tagP.

    \item \textbf{Adiabatic approximation}: bulk modes relax at
    fixed $q$ \tagP.
\end{enumerate}
\end{assumptionbox}

\subsection{What ``Bulk Backreaction'' Means in N2}

\textbf{Physical mechanism:} When the Y-junction node displaces by $q$,
the brane stress-energy distribution $T_{AB}$ changes. The 5D bulk
metric adjusts through the linearized Einstein equations:
\begin{equation}
    \delta G_{AB} = \kappa_5^2\, \delta T_{AB}
    \label{eq:N2_linearized}
\end{equation}
The energy stored in the bulk metric perturbation $h_{AB}(q)$ defines
$V_{\mathrm{bulk}}(q)$.

\textbf{Distinction from N1 and N7:}
\begin{itemize}
    \item N1 (Israel): concerns matching conditions \emph{at} the brane
          surface. Result: tension renormalization only. \textbf{Bounded no-go.}
    \item N7 (thick core): concerns internal energy \emph{inside} the
          junction vertex at scale $\delta$. Result: monotone profiles
          reinforce $V_{\mathrm{geom}}$. \textbf{Bounded insufficiency.}
    \item N2 (bulk): concerns gravitational field energy \emph{in the
          5D bulk} away from the brane. This is the response of the
          bulk geometry to the changed source.
\end{itemize}

\subsection{Scales and Couplings}

\begin{center}
\begin{tabular}{lll}
\toprule
\textbf{Symbol} & \textbf{Meaning} & \textbf{Tag} \\
\midrule
$\kappa_5^2 = 8\pi/\Mfive^3$ & 5D gravitational coupling & [Def] \\
$\Mfive$ & 5D Planck mass & [P] \\
$\sigma$ & Brane tension: $8.82\MeV/\fm^2$ & [Dc] \\
$R$ & Y-junction arm length $\sim 1\fm$ & [Dc] \\
$q$ & Node displacement from Steiner & [Dc] \\
$R_5$ & Compactification radius & [P] \\
$k$ & RS curvature parameter (if applicable) & [P] \\
\bottomrule
\end{tabular}
\end{center}

% ============================================================================
\section{Linearized Bulk Response: Structural Derivation}
\label{sec:P2N2:derivation}
% ============================================================================

\subsection{Total Energy Decomposition}

The total on-shell energy of the static brane-bulk system at junction
displacement $q$ is:
\begin{equation}
    E_{\mathrm{total}}(q) = E_{\mathrm{brane}}(q) + E_{\mathrm{grav}}(q)
    \label{eq:N2_Etotal}
\end{equation}
where $E_{\mathrm{brane}}(q) = V_{\mathrm{geom}}(q) = \tau L_{\mathrm{tot}}(q)$
\tagDc\ is the Nambu--Goto brane energy, and $E_{\mathrm{grav}}(q)$ is the
gravitational energy (bulk + GHY + junction terms evaluated on the
solution to the Einstein equations). The N2 contribution is:
\begin{equation}
    V_{\mathrm{bulk}}(q) \equiv E_{\mathrm{grav}}(q) - E_{\mathrm{grav}}(0)
    \label{eq:N2_Vbulk_def}
\end{equation}

\subsection{Perturbative Expansion}

Let $\bar{g}_{AB}$ be the background metric solving the 5D Einstein
equations with the Steiner brane configuration ($q = 0$) as source.
When the junction displaces to $q > 0$:
\begin{align}
    T_{AB}(q) &= T_{AB}(0) + \delta T_{AB}(q) \label{eq:N2_deltaT} \\
    g_{AB}(q) &= \bar{g}_{AB} + h_{AB}(q) \label{eq:N2_deltag}
\end{align}
where $h_{AB}(q)$ solves the linearized Einstein equations with
source $\delta T_{AB}(q)$.

The gravitational energy at second order in perturbation theory is
(standard result in linearized gravity):
\begin{equation}
    V_{\mathrm{bulk}}(q) = \underbrace{-\int \delta T_{AB}\, \bar{h}^{AB}\,
    d^5x}_{\text{cross-term with background}} \;-\;
    \underbrace{\frac{1}{2}\int \delta T_{AB}\, \delta h^{AB}\,
    d^5x}_{\text{self-energy of perturbation}}
    \label{eq:N2_Vbulk_perturbative}
\end{equation}
where $\bar{h}^{AB}$ is the background metric perturbation (sourced by
$T_{AB}(0)$) and $\delta h^{AB}$ is the perturbation sourced by
$\delta T_{AB}(q)$. \tagDc\ (standard linearized gravity).

\subsection{Cross-Term: Tension Renormalization}
\label{sec:N2:cross}

The first term in Eq.~\eqref{eq:N2_Vbulk_perturbative} involves the
background gravitational field $\bar{h}_{AB}$ evaluated on the source
perturbation $\delta T_{AB}$.

\textbf{Key observation:} The background field $\bar{h}_{AB}$ is
sourced by the Steiner brane configuration $T_{AB}(0)$. For a brane
at fixed $\xi = \xi_0$ in a warped background, the Israel matching
condition gives (Appendix~\ref{app:P2_WP2_Israel}):
\begin{equation}
    \bar{h}_{\mu\nu} \propto A'(\xi_0)\,e^{2A(\xi_0)}\,\eta_{\mu\nu}
    \label{eq:N2_hbar}
\end{equation}
This is \textbf{independent of in-brane position}: the background
gravitational field at $\xi = \xi_0$ is uniform across the brane
plane. \tagDc

Therefore, the cross-term evaluates to:
\begin{equation}
    -\int \delta T_{AB}\, \bar{h}^{AB}\, d^5x
    \;\propto\; A'(\xi_0)\,\sigma \int_{\mathrm{brane}} \delta(\sqrt{h})\, d^4x
    \;\propto\; \delta L_{\mathrm{tot}}(q)
    \;\propto\; V_{\mathrm{geom}}(q) - V_{\mathrm{geom}}(0)
    \label{eq:N2_cross_result}
\end{equation}

\begin{resultbox}[title=Cross-Term Result]
The cross-term in the bulk backreaction is \textbf{proportional to}
$V_{\mathrm{geom}}(q)$. It renormalizes the effective brane tension:
\begin{equation}
    V_{\mathrm{cross}}(q) = -c_{\mathrm{cr}}\,\kappa_5^2\,\sigma\,
    A'(\xi_0)\;V_{\mathrm{geom}}(q)
\end{equation}
where $c_{\mathrm{cr}}$ is a positive dimensionless constant from the
tensor contractions.

\textbf{This does not change the single-well structure.} The renormalized
potential $V_{\mathrm{geom}}^{\mathrm{ren}}(q) = (1 + \Delta\tau/\tau)\,
V_{\mathrm{geom}}(q)$ has the same Steiner minimum and no secondary
extremum.

This result is \textbf{structurally identical} to the N1 finding
(Appendix~\ref{app:P2_WP2_Israel}): the leading-order gravitational
backreaction merely renormalizes the tension. \tagDc
\end{resultbox}

\subsection{Self-Energy Term: The Genuinely New N2 Contribution}
\label{sec:N2:self}

The second term in Eq.~\eqref{eq:N2_Vbulk_perturbative} is the
gravitational self-energy of the source perturbation:
\begin{equation}
    V_{\mathrm{self}}(q) = -\frac{\kappa_5^2}{2}\int\!\!\int
    \delta T_{AB}(x)\; G^{AB,CD}(x,x')\;
    \delta T_{CD}(x')\; d^5x\, d^5x'
    \label{eq:N2_Vself}
\end{equation}
where $G^{AB,CD}(x,x')$ is the linearized gravitational Green's
function in the background $\bar{g}_{AB}$. This is the only term that
could produce structure \emph{different from} $V_{\mathrm{geom}}$.
\tagDc

\subsubsection{Source Structure}

The source perturbation $\delta T_{AB}(q)$ encodes the change in
brane stress-energy when the junction displaces. For a Nambu--Goto
brane:
\begin{equation}
    T^{AB}_{\mathrm{brane}} = -\sigma\, h^{AB}\,
    \frac{\delta^{(1)}(\xi - \xi_0)}{\sqrt{g_{\xi\xi}}}
    \label{eq:N2_Tbrane}
\end{equation}
where $h^{AB}$ is the induced metric on the brane worldsheet.

When the junction node moves from $(0,0,\xi_0)$ to $(q,0,\xi_0)$,
the three arm worldsheets rotate and change length. The perturbation
$\delta T_{AB}$ is the \emph{difference} between the stress-energy
of the displaced and Steiner configurations, localized on the brane
at $\xi = \xi_0$.

\textbf{Key property:} $\delta T_{AB}$ has compact support (it is
nonzero only in the region where the arm worldsheets differ between
$q$ and $0$) and has zero total ``charge'' in the sense that:
\begin{equation}
    \int_{\mathrm{brane}} \delta T_{AB}\, d^3x
    \;\propto\; \sigma\,\delta L_{\mathrm{tot}}(q)
    \label{eq:N2_deltaT_total}
\end{equation}
This is the total mass-energy change, already captured by $V_{\mathrm{geom}}$.
The self-energy contribution depends on the \emph{spatial distribution}
of $\delta T$, not just its total integral.

\subsubsection{Self-Energy Decomposition}

The self-energy~\eqref{eq:N2_Vself} receives contributions from:
\begin{enumerate}
    \item \textbf{Monopole term:} gravitational self-energy of the
    total mass change $\delta M \propto \sigma\,\delta L_{\mathrm{tot}}$.
    This scales as $\delta M^2 \propto (\delta L_{\mathrm{tot}})^2
    \propto q^4/R^2$ at leading order (since $L'_{\mathrm{tot}}(0) = 0$
    by Steiner optimality).

    \item \textbf{Dipole and higher multipole terms:} gravitational
    energy from the \emph{redistribution} of brane mass-energy
    (arms rotating, node shifting). These are the shape-dependent
    contributions.
\end{enumerate}

At leading order in $q$, the dipole contribution dominates because
the Steiner displacement is a shift (dipole-like perturbation):
\begin{equation}
    V_{\mathrm{self}}(q) \approx -\frac{\kappa_5^2\,\sigma^2}{2}
    \sum_{i,j=1}^{3}\int\!\!\int \delta_i(s)\; G_{\mathrm{red}}(r_i(s), r_j(s'))\;
    \delta_j(s')\; ds\, ds'
    \label{eq:N2_Vself_arms}
\end{equation}
where $\delta_i$ encodes the perturbation of arm $i$, $r_i(s)$ is the
arm position, $s$ is the arc-length parameter, and $G_{\mathrm{red}}$
is the Green's function reduced to the brane at $\xi = \xi_0$.

\subsubsection{Leading-Order Scaling}

\begin{derivationbox}[title=Self-Energy Scaling Argument]
The self-energy integral~\eqref{eq:N2_Vself_arms} has the following
dimensional structure:

\textbf{Coupling:} $\kappa_5^2 = 8\pi/\Mfive^3$

\textbf{Source strength:} $\sigma^2$ (brane tension squared)

\textbf{Geometric integral:} involves the reduced Green's function
$G_{\mathrm{red}}$, which for a compact extra dimension of size $R_5$
gives:
\begin{equation}
    G_{\mathrm{red}}(r) = \frac{1}{R_5}\,G_{4D}(r)
    + \sum_{n \geq 1} \frac{2}{R_5}\,
    G_{4D}^{(m_n)}(r)
    \label{eq:N2_Green_KK}
\end{equation}
where $G_{4D}(r) \sim 1/(4\pi^2 r^2)$ is the 4D (massless) scalar
Green's function and $G_{4D}^{(m_n)}$ are the KK-massive contributions
with $m_n = n/R_5$.

\textbf{The zero-mode contribution:}
\begin{equation}
    V_{\mathrm{self}}^{(0)}(q) \sim -\frac{\kappa_5^2\,\sigma^2}{R_5}
    \times R^2 \times \mathcal{F}(q/R)
    \label{eq:N2_Vself_zero}
\end{equation}
where $\mathcal{F}$ is a dimensionless geometric function with
$\mathcal{F}(0) = 0$ and $R^2$ arises from the integration over the
arm extent $\sim R$ and the $1/r^2$ Green's function.

Using $G_4 \equiv G_5/R_5 = 1/(\Mfive^3 R_5)$:
\begin{equation}
    V_{\mathrm{self}}^{(0)}(q) \sim -\frac{\sigma^2\, R^2}{\Mfive^3\, R_5}
    \times \mathcal{F}(q/R)
    \label{eq:N2_Vself_scaling}
    \tagDc
\end{equation}
\end{derivationbox}

\subsection{Leading-Order Area Dependence}
\label{sec:N2:area}

The self-energy of the brane configuration depends on the total brane
area (arm length in the reduced problem). At leading order in the
gravitational coupling:
\begin{equation}
    E_{\mathrm{grav}}(q) = -\frac{c_0\,\kappa_5^2\,\sigma^2}{R_5}\,
    L_{\mathrm{tot}}(q) \times [\text{UV-dependent factor}]
    \label{eq:N2_Egrav_LO}
\end{equation}
where $c_0$ is a positive constant from the tensor contractions and
the UV-dependent factor involves $\log(R/\delta)$ or $R/\delta$
depending on the codimension of the source.

\begin{resultbox}[title=Leading-Order Structure]
The leading-order gravitational self-energy is \textbf{proportional
to the total arm length} $L_{\mathrm{tot}}(q)$, which is proportional
to $V_{\mathrm{geom}}(q)$.

\textbf{Physical reason:} at leading order, the gravitational
self-energy depends on ``how much brane'' there is, not on its shape.
More brane (longer arms at larger $q$) means more gravitational
binding energy (more negative $E_{\mathrm{grav}}$).

\textbf{Consequence:} the leading-order bulk backreaction produces
another tension renormalization:
\begin{equation}
    V_{\mathrm{bulk}}^{(\mathrm{LO})}(q) \propto V_{\mathrm{geom}}(q)
    \label{eq:N2_LO_prop}
\end{equation}
The single-well structure is preserved. \tagDc

This is \textbf{structurally the same result} as N1 (Israel no-go,
Appendix~\ref{app:P2_WP2_Israel}) and is consistent with it:
both the junction matching and the leading bulk self-energy
produce tension renormalization only.
\end{resultbox}

\subsection{Next-to-Leading Order: Shape-Dependent Correction}
\label{sec:N2:NLO}

The shape-dependent correction arises from the \emph{mutual}
gravitational interaction between different parts of the brane:
\begin{equation}
    V_{\mathrm{bulk}}^{(\mathrm{NLO})}(q) = -\frac{\kappa_5^2\,\sigma^2}{2}
    \sum_{i \neq j}\int\!\!\int
    \bigl[\rho_i^{(q)}(s) - \rho_i^{(0)}(s)\bigr]\;
    G_{\mathrm{red}}(r_i^{(q)}(s), r_j^{(q)}(s'))\;
    \bigl[\rho_j^{(q)}(s') - \rho_j^{(0)}(s')\bigr]\;
    ds\, ds'
    \label{eq:N2_VNLO}
\end{equation}
where the sum runs over distinct arm pairs, and the Green's function
evaluates the gravitational interaction between perturbations on
different arms.

\textbf{Sign of NLO term:} The sign depends on the geometry. When
the junction displaces:
\begin{itemize}
    \item Arms~2 and~3 rotate toward each other (their angle narrows
          from $120^\circ$). This brings parts of the brane closer
          together, increasing the mutual gravitational binding
          (more negative contribution).
    \item Arm~1 moves away from arms~2 and~3. This decreases the
          mutual gravitational binding (less negative contribution).
\end{itemize}

The net sign is determined by the competition between these effects.
By the $\Zthree$ symmetry of the Steiner configuration, the leading
$q$-dependent correction starts at $O(q^2)$:
\begin{equation}
    V_{\mathrm{bulk}}^{(\mathrm{NLO})}(q) = -\frac{\kappa_5^2\,\sigma^2\,R}{R_5}
    \times c_{\mathrm{shape}} \times \frac{q^2}{R^2} + O(q^4)
    \label{eq:N2_VNLO_expansion}
\end{equation}
where $c_{\mathrm{shape}}$ is a dimensionless constant whose sign
depends on the detailed arm geometry. \tagDcI

\textbf{Even if $c_{\mathrm{shape}} > 0$} (so the NLO term is
attractive at small $q$, opposing $V_{\mathrm{geom}}''(0) > 0$),
the magnitude is:
\begin{equation}
    \left|\frac{V_{\mathrm{bulk}}^{(\mathrm{NLO})}(q)}{V_{\mathrm{geom}}(q)}\right|
    \sim \frac{\kappa_5^2\,\sigma\,R}{R_5}
    = \frac{8\pi\,\sigma\,R}{\Mfive^3\,R_5}
    \label{eq:N2_ratio}
    \tagDc
\end{equation}

This is the \textbf{$\kappa_5^2$-suppression ratio} that determines
whether the N2 lane can ever produce a relevant contribution.

% ============================================================================
\section{$\kappa_5^2$-Suppression Analysis}
\label{sec:P2N2:suppression}
% ============================================================================

\subsection{Magnitude Estimate}

Using $\sigma = 8.82\MeV/\fm^2$ and $R \approx 1\fm$:
\begin{equation}
    \frac{8\pi\,\sigma\,R}{\Mfive^3\,R_5}
    = \frac{8\pi \times 8.82\MeV}{\Mfive^3 \times R_5 \times \fm}
    \label{eq:N2_suppression_numeric}
\end{equation}

Converting to consistent units ($\hbar = c = 1$, where
$1\fm = 1/(197.3\MeV)$):
\begin{equation}
    8\pi\,\sigma\, R = 8\pi \times 8.82 \times (197.3)^2 \times (1/197.3)\MeV
    = 8\pi \times 8.82 \times 197.3\MeV^2
    \approx 4.35 \times 10^4\MeV^2
    \label{eq:N2_numerator}
\end{equation}

The suppression ratio becomes:
\begin{equation}
    \frac{V_{\mathrm{bulk}}^{(\mathrm{NLO})}}{V_{\mathrm{geom}}}
    \sim \frac{4.35 \times 10^4}{\Mfive^3 \times R_5}
    \label{eq:N2_suppression_ratio}
\end{equation}
where $\Mfive$ is in MeV and $R_5$ in MeV$^{-1}$.

\subsection{$\Mfive$ Threshold for Relevance}

For the bulk backreaction to be competitive with $V_{\mathrm{geom}}$
(ratio $\gtrsim 0.1$), we need:
\begin{equation}
    \Mfive^3 \times R_5 \lesssim 4.35 \times 10^5\MeV^2
    \label{eq:N2_threshold}
\end{equation}

\begin{center}
\begin{tabular}{llll}
\toprule
$\Mfive$ & $R_5 = \pi\bdelta$ & $\Mfive^3 R_5$ & Suppression ratio \\
\midrule
$10\MeV$ & $1.67 \times 10^{-3}\MeV^{-1}$ & $1.67\MeV^2$ & $\sim 2.6 \times 10^4$ \\
$100\MeV$ & $1.67 \times 10^{-3}\MeV^{-1}$ & $1.67 \times 10^3\MeV^2$ & $\sim 26$ \\
$1\;\mathrm{GeV}$ & $1.67 \times 10^{-3}\MeV^{-1}$ & $1.67 \times 10^6\MeV^2$ & $\sim 0.026$ \\
$10\;\mathrm{GeV}$ & $1.67 \times 10^{-3}\MeV^{-1}$ & $1.67 \times 10^9\MeV^2$ & $\sim 2.6 \times 10^{-5}$ \\
$M_{\mathrm{Pl}} \approx 10^{18}\;\mathrm{GeV}$ & $1.67 \times 10^{-3}\MeV^{-1}$ & $\sim 10^{51}\MeV^2$ & $\sim 10^{-47}$ \\
\bottomrule
\end{tabular}
\end{center}

\begin{resultbox}[title=$\kappa_5^2$-Suppression Verdict]
The shape-dependent (NLO) bulk backreaction is competitive with
$V_{\mathrm{geom}}$ \textbf{only if $\Mfive \lesssim 100\MeV$}
(for $R_5 = \pi\bdelta$).

For any $\Mfive$ above $\sim 1\;\mathrm{GeV}$, the bulk backreaction
is suppressed below the $1\%$ level relative to $V_{\mathrm{geom}}$.

For $\Mfive \sim M_{\mathrm{Pl}}$ (standard Randall--Sundrum models):
the suppression is $\sim 10^{-47}$ --- completely negligible.

\textbf{$\Mfive$ is undetermined} \tagP. No derivation on any branch
constrains $\Mfive$ from within EDC. The suppression threshold
$\Mfive \lesssim 100\MeV$ is an unusually low 5D Planck mass ---
far below the electroweak scale. Whether such a low $\Mfive$ is
physically viable within EDC is an open question not addressed here.
\tagOpen
\end{resultbox}

% ============================================================================
\section{Comparison to Accepted Geometric Sector}
\label{sec:P2N2:comparison}
% ============================================================================

\subsection{Structure Summary}

\begin{center}
\begin{tabular}{lllll}
\toprule
\textbf{Contribution} & \textbf{$q$-Dependence} & \textbf{Sign}
& \textbf{Scaling} & \textbf{Creates well?} \\
\midrule
$V_{\mathrm{geom}}$ & $\tau L_{\mathrm{tot}}(q)$ & Repulsive
& $\sigma R$ & NO (single-well) \\
Cross-term (LO) & $\propto V_{\mathrm{geom}}(q)$ & Renormalizes $\tau$
& $\kappa_5^2 \sigma^2 R / R_5$ & NO (same shape) \\
LO self-energy & $\propto L_{\mathrm{tot}}(q)$ & Renormalizes $\tau$
& $\kappa_5^2 \sigma^2 R / R_5$ & NO (same shape) \\
NLO self-energy & shape-dependent $q^2 + \ldots$
& Sign undetermined & $(\kappa_5^2 \sigma R/R_5) \times V_{\mathrm{geom}}$
& Only if $\Mfive \lesssim 100\MeV$ \\
\bottomrule
\end{tabular}
\end{center}

\subsection{Does $V_{\mathrm{bulk}}$ Oppose the Geometric Restoring Force?}

\textbf{At leading order: NO.} The cross-term and LO self-energy are
both proportional to $V_{\mathrm{geom}}(q)$. They renormalize the
tension but preserve the single-well structure.

\textbf{At NLO: POSSIBLY, but suppressed.} The shape-dependent NLO
correction has undetermined sign ($c_{\mathrm{shape}}$ is not
computed here). If $c_{\mathrm{shape}} > 0$, the NLO term could
weaken the Steiner curvature $V''(0)$. But the weakening is
suppressed by the ratio $\kappa_5^2 \sigma R / R_5 \ll 1$ for any
physically reasonable $\Mfive$.

\textbf{Can $V_{\mathrm{bulk}}$ create a secondary minimum?}
For a secondary minimum, the NLO correction would need to
\emph{dominate} $V_{\mathrm{geom}}$ at some $q > 0$. This requires
the suppression ratio to be $\gtrsim 1$, which means
$\Mfive \lesssim O(100\MeV)$ --- an extremely low 5D Planck mass.

% ============================================================================
\section{Structural vs.\ Numeric Closure Firewall}
\label{sec:P2N2:firewall}
% ============================================================================

\textbf{What was structurally derived:}
\begin{itemize}
    \item The decomposition of $V_{\mathrm{bulk}}(q)$ into cross-term
          + LO self-energy + NLO self-energy \tagDc
    \item The cross-term $\propto V_{\mathrm{geom}}(q)$ (tension
          renormalization) \tagDc
    \item The LO self-energy $\propto L_{\mathrm{tot}}(q)$ (tension
          renormalization) \tagDc
    \item The $\kappa_5^2$-suppression scaling of the NLO contribution
          \tagDc
    \item The suppression threshold $\Mfive \lesssim 100\MeV$ \tagDc
\end{itemize}

\textbf{What remains parameterized:}
\begin{itemize}
    \item The 5D Planck mass $\Mfive$ \tagP\ (not derived)
    \item The sign of $c_{\mathrm{shape}}$ (NLO geometric coefficient)
          \tagOpen
    \item The compactification radius $R_5$ \tagP
    \item The background metric $A(\xi)$ \tagP
    \item Whether $\Mfive \lesssim 100\MeV$ is physically viable
          in EDC \tagOpen
\end{itemize}

\textbf{Why no background tuning was used:} The background metrics
(flat and RS-like) were declared \tagP\ before computation
(\S\ref{sec:P2N2:model}). The $\kappa_5^2$-suppression scaling
holds for \emph{both} backgrounds --- it is a structural feature of
linearized 5D gravity, not a consequence of a specific metric choice.

\textbf{Why Variant~2 warped scans are negative prior, not positive
evidence:} The Variant~2 scan (D-N2-3) tested different fixed
background metrics, not junction-sourced perturbations. The present
derivation addresses the perturbative response but confirms the
structural expectation: the response is $\kappa_5^2$-suppressed.

\textbf{Why node-well matching was not used:} $V_B$, $\tau_n$,
$\Delta m_{np}$ appear nowhere in this derivation. The
$\kappa_5^2$-suppression result is independent of these target values.

% ============================================================================
\section{Honest Outcome Classification}
\label{sec:P2N2:outcome}
% ============================================================================

\begin{resultbox}[title=N2 WP2 Outcome: Bounded Insufficiency
    (Structural $\kappa_5^2$-Suppression)]

\textbf{Classification:} Bounded insufficiency within the admissible
N2 model class.

\textbf{What was established:}
\begin{enumerate}
    \item \textbf{Cross-term [Dc]:} The interaction of the source
    perturbation with the background gravitational field is
    \emph{proportional to} $V_{\mathrm{geom}}(q)$. It renormalizes
    the brane tension but preserves the single-well structure.
    Structurally identical to the N1 (Israel) result.

    \item \textbf{Leading-order self-energy [Dc]:} The gravitational
    self-energy of the displaced brane is proportional to the total
    arm length $L_{\mathrm{tot}}(q)$. Another tension renormalization.
    No new structure.

    \item \textbf{NLO shape-dependent correction [Dc]:} The first
    correction with potentially new $q$-structure scales as
    $\kappa_5^2 \sigma R / R_5 \sim \sigma R / (\Mfive^3 R_5)$.
    This is suppressed below the $1\%$ level for any
    $\Mfive \gtrsim 1\;\mathrm{GeV}$.

    \item \textbf{Threshold [Dc]:} The NLO correction can compete
    with $V_{\mathrm{geom}}$ only if $\Mfive \lesssim 100\MeV$,
    an unusually low 5D Planck mass.

    \item \textbf{Sign [OPEN]:} The sign of the NLO shape coefficient
    $c_{\mathrm{shape}}$ is not determined. Even if it opposes
    $V_{\mathrm{geom}}$, the magnitude is suppressed.
\end{enumerate}

\textbf{Why ``bounded insufficiency'' not ``no-go'':}
The $\kappa_5^2$-suppression is parametric, not absolute. If
$\Mfive$ is very low ($\lesssim 100\MeV$), the bulk backreaction
could in principle contribute. But:
\begin{itemize}
    \item $\Mfive$ is undetermined [P] --- no branch derives it.
    \item $\Mfive \lesssim 100\MeV$ is far below any scale
          appearing elsewhere in EDC ($\sigma$, $1/\bdelta$, etc.
          correspond to $\sim\mathrm{GeV}$ scales).
    \item Even with a low $\Mfive$, the NLO sign is undetermined.
\end{itemize}

The lane is not dead (a low-$\Mfive$ model is not excluded within
EDC), but it requires an additional hypothesis ($\Mfive$ very low)
not present in the current donor base.

\textbf{Scope of the insufficiency:}
\begin{itemize}
    \item Linearized 5D Einstein-Hilbert gravity
    \item Flat and RS-warped backgrounds
    \item Nambu--Goto brane arms with distributional source
    \item In-brane node displacement (R3 path)
    \item Any value of $\sigma$, $R$
    \item Parametric in $\Mfive$, $R_5$
\end{itemize}

\textbf{Epistemic tag:} \tagDc\ (structural suppression within
linearized gravity).
\end{resultbox}

\subsection{Impact on Candidate Landscape}

\begin{center}
\begin{tabular}{llll}
\toprule
\textbf{Candidate} & \textbf{Before N2 WP2} & \textbf{After N2 WP2}
& \textbf{Status} \\
\midrule
N1 Israel thin-junction & Bounded no-go & Unchanged & Dead end \\
N7 thick-junction (monotone) & Bounded insufficiency & Unchanged
& Narrowed \\
N2 bulk backreaction & [OPEN] --- primary active & Bounded insufficiency
($\kappa_5^2$-suppressed) & Narrowed \\
\bottomrule
\end{tabular}
\end{center}

% ============================================================================
\section{Summary}
\label{sec:P2N2:summary}
% ============================================================================

\begin{resultbox}[title=N2 WP2 Summary]
\textbf{Question asked:} Can the bulk gravitational backreaction from
a displaced Y-junction generate a non-geometric contribution to $V(q)$?

\textbf{Answer:} The mechanism is real but
\textbf{$\kappa_5^2$-suppressed}: quantitatively irrelevant unless
$\Mfive \lesssim 100\MeV$ [P].

\textbf{Key results:}
\begin{center}
\begin{tabular}{lll}
\toprule
\textbf{Item} & \textbf{Result} & \textbf{Tag} \\
\midrule
Cross-term & $\propto V_{\mathrm{geom}}$ (tension renorm.) & [Dc] \\
LO self-energy & $\propto L_{\mathrm{tot}}$ (tension renorm.) & [Dc] \\
NLO shape correction & Suppressed by $\kappa_5^2 \sigma R / R_5$ & [Dc] \\
Threshold & $\Mfive \lesssim 100\MeV$ for relevance & [Dc] \\
Sign of NLO & Undetermined & [OPEN] \\
$\Mfive$ value & Not derived & [P] \\
\bottomrule
\end{tabular}
\end{center}

\textbf{N2 lane status after WP2:} Narrowed. Bulk backreaction is
$\kappa_5^2$-suppressed for any $\Mfive \gtrsim 1\;\mathrm{GeV}$.
The lane survives only under the additional hypothesis
$\Mfive \lesssim 100\MeV$, which is not derived and has no support
in the current donor base.

\textbf{Cumulative status of all node-well candidates:}
\begin{itemize}
    \item N1 (Israel thin-junction): bounded no-go [Dc]
    \item N7 (thick-junction monotone): bounded insufficiency [Dc]
    \item N2 (bulk backreaction): bounded insufficiency
          ($\kappa_5^2$-suppressed) [Dc]
    \item N5 (Helfrich): falsified [Dc]
    \item N4 ($\xi$-BC): falsified [Dc]
    \item N6 (phenomenological): [P/Cal], no derivation
\end{itemize}

\textbf{No untested candidate remains} within the
$S_{\mathrm{EH}} + S_{\mathrm{NG}}$ action class with
$\Mfive \gtrsim 1\;\mathrm{GeV}$.

\textbf{Anti-smuggling compliance:}
\begin{itemize}
    \item[$\checkmark$] No forbidden donors imported
    \item[$\checkmark$] Background metric declared before computation
    \item[$\checkmark$] $\tau_n$, $V_B$, $\Delta m_{np}$ not used
    \item[$\checkmark$] Variant~2 treated as negative prior, not positive evidence
    \item[$\checkmark$] $\kappa_5^2$-suppression reported honestly
    \item[$\checkmark$] $\Mfive$ not tuned; results parametric
\end{itemize}
\end{resultbox}
